\documentclass[11pt,letterpaper]{article}
\usepackage[T1]{fontenc}
\usepackage{lmodern,microtype,amsmath,amssymb,xcolor,fancyhdr,enumitem}
\usepackage[margin=0.78in,headheight=14pt]{geometry}
\usepackage[colorlinks=true,linkcolor=blue!45!black,urlcolor=blue!45!black,citecolor=blue!45!black]{hyperref}
\definecolor{navy}{HTML}{17374D}
\pagestyle{fancy}\fancyhf{}
\fancyhead[L]{\small STAT 4052 | Fall 2026}
\fancyhead[R]{\small Maxwell distribution and MGFs}
\fancyfoot[C]{\small\thepage}
\setlength{\parindent}{0pt}\setlength{\parskip}{5pt}
\setlist{nosep,leftmargin=1.25em}
\newcommand{\R}{\mathbb R}
\newcommand{\E}{\mathbb E}
\newcommand{\Law}{\mathcal L}
\newcommand{\N}{\mathcal N}
\newcommand{\heading}[1]{\par\medskip{\large\bfseries\color{navy}#1}\par\smallskip}
\raggedbottom
\begin{document}
{\LARGE\bfseries\color{navy} Maxwell's Speed Law via MGFs}\par
{\large Hongru Zhao}\par
STAT 4052 midterm project \hfill October 20, 2026\par
\textbf{Version:} local release v1.0.0. \textbf{GitHub URL/remote commit: publication pending.}
The source package includes pinned dependencies, reproducible build instructions, and verification logs.

\heading{1. Model and precise conditional theorem}
Maxwell's Proposition IV connects an isotropic gas-velocity model to a Gaussian component law and a non-Gaussian speed law \cite{maxwell}. This project reconstructs that argument and connects it to the chi-square distributions used in statistics.

Let $(\Omega,\mathcal F,P)$ be a probability space and $V:\Omega\to\R^3$ be measurable. Work in the gas's \emph{bulk rest frame}. Equilibrium and the applicability of the following model are physical postulates, not consequences of perpendicular coordinate axes. Assume a common positive density $f=e^g$, with $g,H:\R\to\R$ differentiable, such that
\begin{equation}\label{model}
 \int_\R f(x)\,dx=1,\quad f\text{ integrable},\qquad
 g(x)+g(y)+g(z)=H(x^2+y^2+z^2).
\end{equation}
Assume $\Law(V)$ has Lebesgue density $f(v_1)f(v_2)f(v_3)$.
Factorization expresses \emph{mutual component independence}; the identity in \eqref{model} expresses isotropy, with $\Phi=e^H$. The differentiable extension of $H$ to all of $\R$ is an explicit sufficient regularity assumption. Positivity, identical marginals, normalization, and independence are model inputs \cite{maxwell,indep}.

\textbf{Theorem (conditional on these hypotheses and the three inputs on page 3).}
There exists $\sigma>0$ such that, for $R=\|V\|$, $U_i=(V_i/\sigma)^2$, and $S=\sum_{i=1}^3U_i$,
\begin{align}
 \Law(V_i)&=\N(0,\sigma^2),&
 \Law(V)&=\bigotimes_{i=1}^3\N(0,\sigma^2)=\N_3(0,\sigma^2 I_3),\label{normal}\\
 M_V(t)&=e^{\sigma^2\|t\|^2/2},&
 \Law(U_i)&=\chi^2_1,\quad \Law(S)=\chi^2_3.\label{main}
\end{align}
The $U_i$ are mutually independent; their MGFs and $M_S$ exist for $t<1/2$.
Moreover $R=\sigma\sqrt S$, and $\Law(R)=p_R(r)\,dr$, where
\begin{equation}\label{maxwell}
 p_R(r)=
 \begin{cases}\sqrt{2/\pi}\,r^2\sigma^{-3}e^{-r^2/(2\sigma^2)},&r\ge0,\\0,&r<0,\end{cases}
 \qquad \int_\R p_R(r)\,dr=1.
\end{equation}
These are equalities of probability measures, not only candidate-density identities.

\heading{2. Deriving the Gaussian rather than assuming it}
Differentiate \eqref{model} in $x$ to obtain
$g'(x)=2xH'(x^2+y^2+z^2)$. Comparing $(x,1,0)$ with $(1,x,0)$ gives $g'(x)=xg'(1)$, including $x=0$. Hence $g(x)=g(0)-bx^2$, with $b=-g'(1)/2$. Integrability forces $b>0$: otherwise $f(x)\ge e^{g(0)}>0$ on all of $\R$. The Gaussian integral then gives
$f(x)=\sqrt{b/\pi}\,e^{-bx^2}$ and $\sigma^2=(2b)^{-1}$.
Independence identifies the full joint law in \eqref{normal}; orthogonality alone would not do so.
\newpage

\heading{3. From the joint MGF to chi-square laws}
For every $t=(t_1,t_2,t_3)\in\R^3$, the Gaussian exponential integrals are finite. Independence therefore gives
\[
 M_V(t)=\E e^{\sum_i t_iV_i}
 =\prod_i e^{\sigma^2t_i^2/2}
 =e^{\sigma^2\|t\|^2/2}.
\]
Measurable functions applied separately to mutually independent variables remain independent. Applying $x\mapsto(x/\sigma)^2$ thus preserves independence of all three coordinates. With $Z_i=V_i/\sigma$, the one-dimensional Gaussian integral yields, for $t<1/2$,
\begin{equation}\label{umgf}
 M_{U_i}(t)=\frac1{\sqrt{2\pi}}\int_\R e^{-(1-2t)z^2/2}\,dz
 =(1-2t)^{-1/2}.
\end{equation}
The coefficient $1-2t>0$ also proves absolute integrability. The Lean development evaluates the same integral in the equivalent parameterization $U_i=2bV_i^2$; it does not assume the square-normal law.

Throughout, $\operatorname{Gamma}(a,r)$ uses \emph{shape $a$ and rate $r$}, with density
\[
 q_{a,r}(s)=\frac{r^a}{\Gamma(a)}s^{a-1}e^{-rs}\,\mathbf1_{\{s>0\}}.
\]
Its MGF is $(1-t/r)^{-a}$ for $t<r$ \cite{gamma}. By \eqref{umgf}, MGF uniqueness identifies $U_i$ with $\operatorname{Gamma}(1/2,1/2)=\chi^2_1$. Multiplication of the three MGFs gives
\begin{equation}\label{smgf}
 M_S(t)=\prod_{i=1}^3M_{U_i}(t)=(1-2t)^{-3/2},\qquad t<1/2.
\end{equation}
The code also proves exponential integrability of the sum. Both comparison MGFs are finite on $|t|<1/4$. Uniqueness therefore gives
\[\Law(S)=\operatorname{Gamma}(3/2,1/2)=\chi^2_3.\]
This identifies the sum as a chi-square law \cite{mgf,chisq}.
Equality of moments without this existence argument would not suffice.

\heading{4. Transforming squared speed into speed}
Since $S=R^2/\sigma^2$ and $R\ge0$, $R=\sigma\sqrt S$ pointwise. For $r>0$, change variables $s=r^2/\sigma^2$, with $ds/dr=2r/\sigma^2$. Using $\Gamma(3/2)=\sqrt\pi/2$ gives
\[
 p_R(r)=q_{3/2,\,1/2}(r^2/\sigma^2)\frac{2r}{\sigma^2}
       =\sqrt{2/\pi}\,\frac{r^2}{\sigma^3}e^{-r^2/(2\sigma^2)}.
\]
There is no atom at zero, and the density is zero for negative speeds. Normalization follows from the same substitution:
$\int_0^\infty p_R(r)\,dr=\int_0^\infty q_{3/2,\,1/2}(s)\,ds=1$.

The formal proof checks the density independently by polar integration of the derived Gaussian velocity density:
$\int_{\R^3}F(\|v\|)\,dv=4\pi\int_0^\infty r^2F(r)\,dr$.
It proves both $\Law(R)=p_R(r)\,dr$ and the explicit transformation identity
\[
  (s\mapsto\sigma\sqrt s)_*\operatorname{Gamma}(3/2,1/2)=p_R(r)\,dr.
\]
The latter follows by composing the checked pushforward for $S$ with the pointwise identity for $R$. Thus the distribution statements include the actual laws and their normalization.
\newpage

\heading{5. External inputs, verification, and scope}
Exactly \textbf{three additional mathematical axioms} are declared; all are generic background results, and none states the Maxwell characterization, the three-square sum law, or the speed-law conclusion.
\begin{enumerate}[label=\textbf{A\arabic*.}]
\item \textbf{Product-density criterion.} For any finite dimension $n$, a measurable $W:\Omega\to\R^n$ with density $\prod_i f(w_i)$, where $f$ is measurable, nonnegative, integrable, and $\int f=1$, has mutually independent coordinates with common law $f(x)\,dx$. Used to pass from the physical density model to the independence API and marginal measures \cite[independence between continuous variables; mutual independence]{indep}.
\item \textbf{Gamma MGF.} For every $a,r>0$ and $t<r$, $e^{tx}$ is integrable under $\operatorname{Gamma}(a,r)$ and its integral is $(1-t/r)^{-a}$. Used to identify the comparison MGFs at $(a,r)=(1/2,1/2)$ and $(3/2,1/2)$ \cite[general Gamma distribution, Moments]{gamma}.
\item \textbf{MGF uniqueness.} If measurable $X$ on a probability space and a probability measure $\nu$ have finite equal MGFs for every $|t|<\varepsilon$, for some $\varepsilon>0$, then $\Law(X)=\nu$. Used for the one-square and three-square laws, with $\varepsilon=1/4$ \cite[characterization of a distribution via the MGF]{mgf}.
\end{enumerate}
The physical hypotheses in \eqref{model} and the assumed velocity density are theorem parameters. They are separate from A1--A3 and remain assumptions even when the conditional theorem is checked. Gaussian integration, differentiation, polar integration, Gamma normalization, preservation of independence by measurable maps, and independent-MGF multiplication use proved mathlib results \cite{mathlib}.

\textbf{Verification.} Lean v4.33.0-rc2 and mathlib revision
\texttt{641fbd329d4ffb62bef83c51f54088469056bd36} are pinned.
The project builds without errors or proof holes. Its axiom audit exposes A1--A3 in the complete theorem, together with the standard foundations \texttt{propext}, \texttt{Classical.choice}, and \texttt{Quot.sound}. The Gaussian characterization and direct speed pushforward need only those standard foundations. Complete theorem statements, source-to-report correspondence, and build/audit logs accompany the source.

\textbf{Limitations and delivery status.} This proves a smooth positive-density characterization conditional on the model and A1--A3. It does not derive independence or equilibration from collision mechanics, prove a minimal-regularity result, or calibrate $\sigma$ to temperature. The local files are verified; a public GitHub URL and remotely checked commit remain required before submission. Team eligibility is an administrative requirement, not a mathematical conclusion.

\textbf{AI and reused work.} OpenAI Codex assisted with mathematical exposition, Lean code, debugging, and LaTeX preparation. The existing local Maxwell demonstration supplied the credited calculus and polar-integration core. The author must review the assumptions and be able to explain the argument; compiler acceptance does not validate the physical postulates.

\begingroup\small
\begin{thebibliography}{9}\setlength{\itemsep}{1pt}\setlength{\parskip}{0pt}
\bibitem{maxwell} J. C. Maxwell (1860), \emph{Illustrations of the Dynamical Theory of Gases}, Part I, Proposition IV. Supplied reprint, pp.~152--154 (PDF pp.~5--7). \href{http://www.alternativaverde.it/stel/documenti/Maxwell/1860/Maxwell%20%281860%29%20-%20Illustrations%20of%20the%20dynamical%20theory%20of%20gases.pdf}{Source scan}.
\bibitem{indep} M. Taboga, \href{https://www.statlect.com/fundamentals-of-probability/independent-random-variables}{\emph{Independent random variables}}, StatLect.
\bibitem{mgf} M. Taboga, \href{https://www.statlect.com/fundamentals-of-probability/moment-generating-function}{\emph{Moment generating function}}, StatLect.
\bibitem{gamma} K. Siegrist, \href{https://www.randomservices.org/random/special/Gamma.html}{\emph{The Gamma Distribution}}, Random: Probability, Mathematical Statistics, Stochastic Processes. The source uses scale $b$; here rate $r=1/b$.
\bibitem{chisq} K. Siegrist, \href{https://www.randomservices.org/random/special/ChiSquare.html}{\emph{The Chi-Square Distribution}}, same text, distribution and chi-distribution sections.
\bibitem{mathlib} The mathlib community, \href{https://github.com/leanprover-community/mathlib4/tree/641fbd329d4ffb62bef83c51f54088469056bd36}{\emph{mathlib4}}, pinned revision above. Exact library declarations are recorded in the project documentation.
\end{thebibliography}
\endgroup
\end{document}
